% Generated from the Lean dependency data; see the comment at the top of content.tex.
\chapter{General spin groups and GGP triples}
\label{chap:gspin}
This chapter follows Thesis Ch.~5. Throughout, $F$ is a $p$-adic field (a finite extension of
$\mathbb{Q}_p$) and $\psi$ a nontrivial continuous additive character of $F$. Quadratic spaces are
nondegenerate and finite-dimensional; $b$ is the bilinear form with $q(v)=b(v,v)$.

\section{Quadratic spaces and the general spin group}
\begin{definition}[The general spin group]
\label{GSpin}
\lean{GSpinGGP.GSpin}
\leanok
For a quadratic space $(V,q)$ with Clifford algebra $C(V)$ and even part $C^+(V)$, the general spin
group is
\[ \GSpin(V) = \{ g \in C^+(V)^\times : g V g^{-1} = V \} \]
(Thesis \S5.1.2). Conjugation gives a surjection $P:\GSpin(V)\to \SO(V)$ on $\Fbar$-points with
kernel the scalars $\mathbb{G}_m$. (This is not the spin group of unit spinor norm.)
\end{definition}

\begin{definition}[The Lie algebra $\mathfrak{gspin}(V)$]
\label{gspin}
\lean{GSpinGGP.gspin}
\leanok
\uses{GSpin}
$\gspin(V) = \{ X \in C^+(V) : [X,v] \in V \text{ for all } v\in V\}$, a Lie subalgebra of
$C(V)$ for the commutator bracket (Thesis Prop.~5.1.3.1); $\gspin(V)=\mathfrak{so}(V)\oplus F$.
\end{definition}

\begin{theorem}[Witt decomposition]
\label{witt_decomposition}
\lean{GSpinGGP.witt_decomposition}
\leanok
A quadratic space is $V \cong H^{w}\perp V_{\mathrm{an}}$ with $H$ the hyperbolic plane and
$V_{\mathrm{an}}$ anisotropic, where $w=w(V)$ is the Witt index (the maximal dimension of a totally
isotropic subspace). In particular $\dim V_{\mathrm{an}} = \dim V - 2w(V)$, so $\GSpin(V)$ is
quasi-split iff $\dim V - 2w(V)\le 2$.
\end{theorem}
\begin{proof}
\leanok
Proved from Tau Ceti's Witt decomposition of a regular form (\cite{Lam}, Ch.~I \S4).
\end{proof}

\begin{theorem}[Witt cancellation]
\label{witt_cancellation}
\lean{GSpinGGP.witt_cancellation}
\leanok
Over a field of characteristic $\neq 2$, if $U\perp X \cong U\perp Y$ then $X\cong Y$
(all forms finite-dimensional, possibly degenerate).
\end{theorem}
\begin{proof}
\leanok
Split off the radicals and apply Tau Ceti's Witt cancellation for a regular summand (\cite{Lam}, Ch.~I \S4).
\end{proof}

\begin{theorem}[Quadratic forms over a $p$-adic field]
\label{padic_quadratic_classes}
\lean{GSpinGGP.padic_quadratic_classes}
\leanok
Let $Q$ be a nondegenerate quadratic form on $F^m$. The isometry classes of nondegenerate forms
on $F^m$ with the same discriminant as $Q$ are finite in number, $n\le 2$ of them, and admit a system of
representatives containing $Q$; $n=2$ if $m\ge 3$, $n=1$ if $m\le 1$, and for $m=2$, $n=1$ iff
$-\det Q$ is a square.
\end{theorem}
\begin{proof}
\leanok
Proved from Tau Ceti's classification of regular forms over a nonarchimedean local field by rank, discriminant and Hasse invariant, and its list of the realizable invariants (\cite{Lam}, Ch.~VI \S2; O'Meara \S63), through the fact that a $p$-adic field is a nonarchimedean local field.
\end{proof}

\begin{theorem}[Cartan--Dieudonn\'e, $\SO$ form]
\label{cartan_dieudonne_even}
\lean{GSpinGGP.cartan_dieudonne_even}
\leanok
In a nondegenerate quadratic space over a field of characteristic $\neq 2$, every isometry of
determinant $1$ is a product $\prod_i s_{u_i}s_{w_i}$ of pairs of reflections in anisotropic vectors.
\end{theorem}
\begin{proof}
\leanok
Proved from Tau Ceti's factorization of an element of $\SO$ into an even number of reflections (\cite{Artin}, Ch.~III; \cite{Lam}, Ch.~I), grouped in pairs.
\end{proof}

\begin{literature}[Centralizers of semisimple elements of $\GSpin$]
\label{LIT-steinberg_centralizer_connected}
\lean{GSpinGGP.steinberg_centralizer_connected}
\leanok
\uses{GSpin}
Let $x\in\GSpin(V)$ be semisimple, $\bar x=P(x)$, and $W_{\pm1}=\ker(\bar x\mp 1)$. Then
$Z_{\GSpin(V)}(x)=P^{-1}(Z_{\SO(V)}(\bar x)^0)$; explicitly, $g$ commutes with $x$ iff $P(g)$
commutes with $\bar x$ and $\det P(g)|_{W_{+1}} = \det P(g)|_{W_{-1}} = 1$. (Characteristic $0$.)

\assumedfrom{\cite{Steinberg}, Thm.~8.1 (the derived group $\Spin(V)$ is simply connected)}
\end{literature}

\begin{literature}[Regular nilpotent orbits and quasi-splitness]
\label{LIT-nilReg_nonempty_iff_quasiSplit}
\lean{GSpinGGP.nilReg_nonempty_iff_quasiSplit}
\leanok
\uses{GSpin, gspin}
$\Nilreg(\gspin(V))\neq\emptyset$ iff $\GSpin(V)$ is quasi-split, i.e.\ $\dim V-2w(V)\le 2$.

\assumedfrom{\cite{Wal01}, Ch.~I; \cite{Wal10}, \S7.1; \cite{Luo}, \S2.1; Thesis \S5.1.4}
\end{literature}

\section{GGP triples}
\begin{definition}[Admissible pairs]
\label{GGPData}
\lean{GSpinGGP.GGPData}
\leanok
A GGP datum is a quadratic space $W$ on $F^m$ together with $r\ge 0$ and $\nu\in F^\times$. Put
$Z$ the split space of dimension $2r+1$ with basis $z_{-r},\dots,z_r$, $b(z_i,z_j)=\nu\delta_{i,-j}$,
and $V = W\oplus^\perp Z$ (Thesis \S5.2).
\end{definition}

\begin{definition}[The group $G$]
\label{G}
\lean{GSpinGGP.GGPData.G}
\leanok
\uses{GGPData, GSpin}
$G = \GSpin(W)\times\GSpin(V)$ (Thesis \S5.2), as a subgroup of $(C(W)\times C(V))^\times$.
Its Lie algebra is $\mathfrak g = \gspin(W)\oplus\gspin(V)$, and $A_G$ denotes its split centre.
\end{definition}

\begin{definition}[The unipotent subgroup $N$]
\label{N}
\lean{GSpinGGP.GGPData.N}
\leanok
\uses{GGPData, GSpin}
$N\subset \GSpin(V)$ is the unipotent radical of the stabilizer $P_V$ in $\GSpin(V)$ of the isotropic
flag $\langle z_r\rangle\subset\langle z_r,z_{r-1}\rangle\subset\dots\subset\langle z_r,\dots,z_1\rangle$
(Thesis \S5.2).
\end{definition}

\begin{definition}[The subgroup $H$]
\label{H}
\lean{GSpinGGP.GGPData.H}
\leanok
\uses{G, GGPData, GSpin, N}
$H = \GSpin(W)\ltimes N\subset G$: the pairs $(a, an)$ with $a\in\GSpin(W)$, $n\in N$
(Thesis \S5.2). $Z_{G,H}=A_{G,H}$ is the diagonal central $\mathbb{G}_m$.
\end{definition}

\begin{definition}[The character $\xi$]
\label{xi}
\lean{GSpinGGP.GGPData.xi}
\leanok
\uses{G, GGPData, GSpin, H}
$\xi(h) = \psi(\lambda(h))$, where $\lambda(an) = \sum_{i=0}^{r-1} b(z_{-i-1}, n z_i)$ is a
homomorphism $H\to\mathbb{G}_a$ (Thesis \S5.2).
\end{definition}

\begin{definition}[Quasi-split triples]
\label{IsQuasiSplit}
\lean{GSpinGGP.GGPData.IsQuasiSplit}
\leanok
\uses{GGPData}
$G$ is quasi-split iff both $W$ and $V$ are, i.e.\ $\dim U-2w(U)\le 2$ for $U=W,V$ (Thesis \S6.4).
\end{definition}

\begin{definition}[Central characters]
\label{CentralCharOn}
\lean{GSpinGGP.CentralCharOn}
\leanok
\uses{G, GGPData, GSpin}
A continuous unitary character $\chi: A_G(F)\to\mathbb{S}^1$ trivial on $A_{G,H}(F)$
(Thesis \S7.1). It is fixed throughout the trace formula.
\end{definition}

\begin{definition}[Whittaker data]
\label{WhittakerDatum}
\lean{GSpinGGP.WhittakerDatum}
\leanok
\uses{G, GGPData, GSpin}
A Whittaker datum of $G$ (Thesis \S4.4): a Borel subgroup $B$ (stabilizer of full maximal isotropic
flags of $W$ and $V$) with unipotent radical $U$, and a generic continuous unitary character $\chi$
of $U(F)$, i.e.\ its stabilizer in $B(F)$ is $Z_G(F)U(F)$. (The Thesis takes $G(F)$-conjugacy classes
of such pairs.)
\end{definition}

\begin{literature}[Luo's orbit $\mathcal O_\nu$]
\label{LIT-exists_regularNil_inOrbitNu}
\lean{GSpinGGP.exists_regularNil_inOrbitNu}
\leanok
\uses{GGPData, IsQuasiSplit, gspin}
Let $G$ be quasi-split. If $\dim V$ is even and $\ge 4$, some regular nilpotent $X\in\gspin(V)$
has image in $\mathfrak{so}(V)$ in the orbit $\mathcal O_\nu$ (it kills an anisotropic line $Fd_0$ and is
one Jordan block on $U_1=d_0^\perp\cong\langle\nu\rangle\oplus H^k$). If $\dim W$ is even and $\ge 4$,
some regular nilpotent $X\in\gspin(W)$ has image in $\mathcal O_{-\nu}$.

\assumedfrom{\cite{Wal10}, \S7.1; \cite{Luo}, \S2.1, \S7.1; \cite{Wal01}, Ch.~I; Thesis \S5.1.4}
\end{literature}

\begin{literature}[Germs vanish off quasi-split centralizers]
\label{LIT-nilReg_centralizer_nonempty_imp}
\lean{GSpinGGP.nilReg_centralizer_nonempty_imp}
\leanok
\uses{G, GGPData, GSpin, gspin}
For semisimple $x\in\GSpin(W)(F)$ embedded diagonally in $G$, with $W'_x=\ker(P(x)-1)$ and
$V'_x = W'_x\oplus Z$: if $\Nilreg(\mathfrak g_x)\ne\emptyset$ then $W'_x$ and $V'_x$ are quasi-split.

\assumedfrom{\cite{Wal01}, Ch.~I applied to the factors of $G_x$; \cite{BP}, Prop.~4.5.1(i)}
\end{literature}

\begin{literature}[Regular nilpotent orbits of GGP centralizers]
\label{LIT-nilReg_centralizer_nonempty_iff}
\lean{GSpinGGP.nilReg_centralizer_nonempty_iff}
\leanok
\uses{G, GGPData, GSpin, gspin}
For semisimple $x\in\GSpin(W)(F)$ whose moved part $G''_x$ is a torus (true for every $x$ contributing
to $\Gamma(G,H)$): $\Nilreg(\mathfrak g_x)\neq\emptyset$ iff $W'_x$ and $V'_x=W'_x\oplus Z$ are
quasi-split.

\assumedfrom{\cite{Wal01}, Ch.~I; \cite{Wal10}, \S7.1}
\end{literature}

\begin{literature}[Multiplicity at most one]
\label{LIT-emoryTakeda_mult_le_one}
\lean{GSpinGGP.emoryTakeda_mult_le_one}
\leanok
\uses{G, GGPData, GSpin, H, IsTempered, UnitaryRep, xi}
For every GGP triple and every irreducible tempered unitary representation $\pi$ of $G(F)$,
$m(\pi)=\dim\Hom_{H(F)}(\pi^\infty,\xi)\le 1$ (Thesis Thm.~5.5.1).

\assumedfrom{\cite{ET}, Thm.~1.3 ($r=0$), with the reduction of \cite{BP}, \S6.3 and \cite{GGP}, \S15; directly \cite{Yan}}
\end{literature}

